% Part: proof-theory
% Chapter: natural-deduction
% Section: rules-proofs

\documentclass[../../../include/open-logic-section]{subfiles}

\begin{document}

\olfileid{pt}{nat}{rp}

\olsection{قاعدې او \usetoken{P}{proof}}

لومړے ځينې تعريفونه ورکوو او ځينې اصطلاحات ټاکو۔

د \Log{N1c} او~\Log{N1i} هغه دوه قاعدې وګورئ چې د $!A
\lif !B$ يادونه کوي۔
\[
\bottomAlignProof
\AxiomC{$\Discharge{!A}{x}$}
\shortDeduce
\DeduceC{$!B$}
\DischargeRule{\Intro{\lif}}{x}
\UnaryInfC{$!A \lif !B$}
\DisplayProof
\qquad
\bottomAlignProof
\AxiomC{$!A \lif !B$}
\AxiomC{$!A$}
\RightLabel{\Elim{\lif}}
\BinaryInfC{$!B$}
\DisplayProof
\]
د \Elim{\lif} قاعده آسانه ده۔ له کرښې بره !!{formula}s
\emph{مقدمې} دي۔ کيڼ اړخ ته !!{formula} هماغه \emph{اصلي}
!!{formula}~$!A \lif !B$ دے۔ ټولې د !!{elimination} قاعدې
يوه داسې مقدمه لري چې تل په تر ټولو کيڼ اړخ کښې وي۔ دې ته د
استنتاج \emph{لويه} مقدمه وايي۔ که نورې مقدمې هم وي، لکه په دې
حالت کښې، هغو ته \emph{کوچنۍ} مقدمې وايي۔ د !!{introduction}
قاعدو مقدمه يا مقدمې هم لويې مقدمې بلل کېږي۔ له کرښې لاندې
!!{formula} \emph{پايله} ده۔

د \Intro{\lif} قاعده يوازې يوه مقدمه لري، او دلته پايله
اصلي فارمول~$!A \lif !B$ دے۔ د !!{introduction} ټولې قاعدې
همداسې دي: پايله اصلي فارمول دے؛ د هغه لويه عملګره هماغه ده چې
``داخلېږي''۔

د $\Discharge{!A}{x}$ نښه دا معنا لري۔ په !!a{proof} کښې
د \Intro{\lif} استنتاج !!a{formula}~$!B$ د مقدمې په توګه اخلي۔
په~$!B$ ختمېدونکے فرعي !!{proof} يو شمېر پيلني !!{formula}s
لري۔ دې ته \emph{فرضونه} وايي۔ !!^a{proof} چې له يوه يا ډېرو
د~$!A$ په بڼه فرضونو څخه~$!B$ ثابتوي، !!a{proof} دے چې $!A$
څخه~$!B$ نتيجه کېږي۔ چې د \Intro{\lif} قاعده وکاروو،~$!A
\lif !B$ نتيجه اخلو او پايله نور په فرض~$!A$ پورې نۀ تړل کېږي۔
د دې ښودلو دپاره، په هغه !!{proof} کښې چې د~\Intro{\lif}
استنتاج هم لري، د~$!A$ په بڼه فرضونه \emph{ايستل کېږي} يا
\emph{بندېږي}۔ قاعده ټاکي چې کوم فرضونه ايستل کېدے شي۔ د
\Intro{\lif} په حالت کښې چې پايله ئې $!A \lif !B$ وي، دا
د~$!A$ په بڼه فرضونه دي، يعنې د شرطي فارمول له مخکينۍ برخې سره
يو شان دي۔ کله کله بايد وښيو چې په کومه قاعده کوم فرضونه ايستل
کېږي؛ د فرض ايستلو نښه~$x$ همدا ښيي۔ د فرضونو ايستونکې قاعدې
\emph{اړ نۀ دي} چې د اړينې بڼې \emph{ټول} فرضونه يا ان يو فرض
هم وباسي۔ مثلاً، دا سم دے، که څه هم لومړے \Intro{\lif}
هېڅ فرض نۀ باسي:
\[
\AxiomC{$\Discharge{!A}{y}$}
\DischargeRule{\Intro{\lif}}{x}
\UnaryInfC{$!B \lif !A$}
\DischargeRule{\Intro{\lif}}{y}
\UnaryInfC{$!A \lif (!B \lif !A)$}
\DisplayProof
\]
چې قاعده هېڅ فرض نۀ باسي، د فرض ايستلو نښه پرېښودے شو؛ په دې
رسم کښې لومړۍ نښه پرېښودے شي۔
\textbf{سرچينه‌يي سمون OLCMP-144:} په اصلي رسم کښې نښه شته،
که څه هم اړوند فرض نشته۔ د نښې پرېښودل اختيار دے؛ د سرچينې
``لکه دلته مو پرېښې ده'' خبره له رسم سره سمه نۀ ده۔

په دې !!{proof} کښې
\[
\AxiomC{$\Discharge{!A}{x}$}
\AxiomC{$\Discharge{!A}{y}$}
\RightLabel{\Intro{\land}}
\BinaryInfC{$!A \land !A$}
\RightLabel{\Elim{\land}}
\UnaryInfC{$!A$}
\DischargeRule{\Intro{\lif}}{x}
\UnaryInfC{$!A \lif !A$}
\DischargeRule{\Intro{\lif}}{y}
\UnaryInfC{$!A \lif (!A \lif !A)$}
\DisplayProof
\]
لومړے \Intro{\lif} استنتاج کيڼ فرض $!A^x$ باسي او دويم
ښي فرض~$!A^y$ باسي۔ يعنې ټول په~$x$ نښه شوي فرضونه هغه
\Intro{\lif} باسي چې نښه ئې~$x$ وي، او ټول په~$y$ نښه شوي
فرضونه هغه باسي چې نښه ئې~$y$ وي۔ د دې د سم کار دپاره بايد
ټول په يوه نښه شوي فرضونه هماغه يو !!{formula} وي۔

د کميت قاعدې لږې پېچلې دي۔ مثلاً، په \Log{N1i} کښې د
$\lexists$ قاعدې دا دي:
\[
\bottomAlignProof
\AxiomC{$!A(t)$}
\RightLabel{\Intro{\lexists}}
\UnaryInfC{$\lexists[x][!A(x)]$}
\DisplayProof
\qquad
\bottomAlignProof
\AxiomC{$\lexists[x][!A(x)]$} 
\AxiomC{$\Discharge{!A(c)}{x}$}
\shortDeduce
\DeduceC{$!B$}
\DischargeRule{\Elim{\lexists}}{x}
\BinaryInfC{$!B$} \DisplayProof 
\bottomAlignProof
\] په
\Intro{\lexists} قاعده کښې مقدمه~$!A(t)$ ده، چې $t$~يو
تړلے ترم دے، يعنې داسې ترم چې متغيرونه نۀ لري۔ دا استنتاج سم
دے، يعنې د \Intro{\lexists} له شېمايي قاعدې سره سمون لري، که
يو !!{formula}~$!A$، يو متغير~$x$ او يو تړلے ترم~$t$ داسې وي
چې پايله $\lexists[x][!A]$ او مقدمه~$\Subst{!A}{t}{x}$ وي۔
د \Elim{\lexists} قاعده د~$!A(c)$ په بڼه فرضونه ايستلو ته
اجازت راکوي؛ يعنې د هر استنتاج دپاره يو !!{constant}~$c$ وي او
د $\Subst{!A}{c}{x}$ په بڼه فرضونه ايستل کېدے شي۔ د يو
استنتاج ټول ايستل شوي فرضونه بايد هماغه~$c$ وکاروي۔ دې ثابت
ته د \Elim{\lexists} استنتاج \emph{ځانګړے متغير} وايي۔
استنتاج يوازې هغه وخت سم دے چې $c$~له ايستل شوي~$!A(c)$ پرته
په بل پرانيستي فرض کښې نۀ راځي، په~$!B$ کښې نۀ راځي او
په~$!A$ کښې نۀ راځي۔ دې ته \emph{د ځانګړي متغير شرط} وايي۔

په \Log{N1c} او~\Log{N1i} کښې !!^{proof}s د فارمولونو
متناهي ونې دي چې له نښه لرونکو فرضونو څخه په استقرايي ډول د
استنتاج د قواعدو په وسيله جوړېږي۔ هره پاڼه يو فرض دے، چې کېدے
شي د فرض ايستلو نښه ولري، او هر بل !!{formula} د نظام د يوې
استنتاجي قاعدې په وسيله له هغو !!{formula}s څخه نتيجه کېږي چې
نېغ ئې بره وي۔ که له استنتاج بره په ونه کښې د فرض په توګه د
!!{formula} يو راتګ له هغه بره استنتاج نۀ وي ايستلے، هغه
په دغه استنتاج کښې \emph{پرانيستے} بلل کېږي۔
\textbf{سرچينه‌يي سمون OLCMP-140:} د N1 ونې د فارمولونو دي،
نۀ د سېکوېنټونو؛ پيلني توکي ئې فرضونه دي۔ ورپسې اصلي تعريف او
د N1 جدول همدغه جوړښت ښيي۔

\begin{defn}\ollabel{defn:proof}
په \Log{N1c} او~\Log{N1i} کښې \emph{!!^{proof}s} د
!!{formula}s متناهي ونې دي چې په استقرايي ډول داسې جوړېږي:
\begin{enumerate}
  \item\ollabel{defn:prf:assumption} هر !!{formula} چې د فرض
  ايستلو نښه لکه $x$ يا $y$ ولري، په ځانله ډول !!a{proof} دے۔
  په هغه !!a{proof} کښې چې يوازې يو !!{formula} لري، دغه
  !!{formula} پرانيستے فرض ګڼل کېږي۔
  \item\ollabel{defn:prf:one-prem} که $R$ يوه قاعده وي چې يوه،
  دوه يا درې مقدمې $!A_1$، $!A_2$، $!A_3$ او پايله~$!A$
  ولري، او $\delta_1$، $\delta_2$، $\delta_3$ هغه !!{proof}s
  وي چې په ترتيب په $!A_1$، $!A_2$، $!A_3$ ختمېږي، نو
  په ترتيب دا
  \[
  \AxiomC{}
  \RightLabel{$\delta_1$}
  \DeduceC{$!A_1$}
  \RightLabel{$R$}
  \UnaryInfC{$!A$}
  \DisplayProof
\qquad
  \AxiomC{}
  \RightLabel{$\delta_1$}
  \DeduceC{$!A_1$}
  \AxiomC{}
  \RightLabel{$\delta_2$}
  \DeduceC{$!A_2$}
  \RightLabel{$R$}
  \BinaryInfC{$!A$}
  \DisplayProof
\qquad
  \AxiomC{}
  \RightLabel{$\delta_1$}
  \DeduceC{$!A_1$}
  \AxiomC{}
  \RightLabel{$\delta_2$}
  \DeduceC{$!A_2$}
  \AxiomC{}
  \RightLabel{$\delta_3$}
  \DeduceC{$!A_3$}
  \RightLabel{$R$}
  \TrinaryInfC{$!A$}
  \DisplayProof
  \]
  !!a{proof} دے، په دې شرط چې د !!{proof}s د فرضونو نښې
  سازګارې وي، يعنې دوه بېل !!{formula}s د فرض ايستلو يوه نښه
  و نۀ لري، او استنتاجونه د~$R$ سمې کارونې وي۔
  \textbf{سرچينه‌يي سمون OLCMP-141:} درېيم فرعي ثبوت په
  خپلې درېيمې مقدمې ختمېږي؛ د رسم نقل شوے دويم اندېکس سم شوے دے۔
  \item په نوي !!{proof} کښې د تر ټولو بره !!{formula} راتګونه
  د !!{proof} \emph{پرانيستي فرضونه} دي، که د $\delta_1$،
  $\delta_2$ يا~$\delta_3$ پرانيستي فرضونه وي، پرته له هغو
  راتګونو چې قاعده ئې باسي۔ که د قاعدې نښه~$x$ يا $x\, y$
  وي، نو د \Intro{\lif} او~\Intro{\lnot} دپاره دا په
  $\delta_1$ کښې ټول په~$x$ نښه شوي پرانيستي فرضونه دي؛ د
  \Elim{\lexists} دپاره په~$\delta_2$ کښې ټول په~$x$ نښه
  شوي پرانيستي فرضونه دي؛ او د \Elim{\lor} دپاره په
  $\delta_2$ کښې ټول په~$x$ نښه شوي او په~$\delta_3$ کښې
  ټول په~$y$ نښه شوي پرانيستي فرضونه دي۔
\end{enumerate}
\end{defn}

هغه !!{proof}s چې د !!a{proof} په استقرايي جوړښت کښې
کارول کېږي، د هغه \emph{فرعي !!{proof}s} بلل کېږي۔ هغه فرعي
!!{proof}s چې د يو استنتاج په مقدمو ختمېږي، د هغه استنتاج
\emph{سمدستي فرعي !!{proof}s} بلل کېږي۔ د \Log{N1c} او
\Log{N1i} قاعدې په \olref{tab:N1} کښې ورکړل شوې دي۔
\textbf{سرچينه‌يي سمون OLCMP-142:} دلته د اصلي N1 جدول او
ورپسې واردات دي؛ د بل نظام نقل شوے نوم سم شوے دے۔

\subfile{rules-N1}

\begin{defn}
د !!a{proof}~$\delta$ \emph{!!{height}}~$\pheight{\delta}$
په استقرايي ډول داسې تعريفېږي۔
\begin{enumerate}
  \item که $\delta$ يوازې يو فرض ولري، $\pheight{\delta} = 0$۔
  \item که $\delta$ په يوه مقدمه لرونکي استنتاج ختمېږي چې
  سمدستي فرعي !!{proof} ئې~$\delta_1$ وي، نو
  $\pheight{\delta} = \pheight{\delta_1} + 1$۔
  \item که $\delta$ په دوو مقدمو لرونکي استنتاج ختمېږي چې
  سمدستي فرعي !!{proof}s ئې~$\delta_1$ او~$\delta_2$ وي، نو
  $\pheight{\delta} =
  \max(\pheight{\delta_1}, \pheight{\delta_2}) + 1$۔
  \item که $\delta$ په \Elim{\lor} ختمېږي چې سمدستي فرعي
  !!{proof}s ئې~$\delta_1$، $\delta_2$ او~$\delta_3$ وي، نو
  $\pheight{\delta} =
  \max(\pheight{\delta_1}, \pheight{\delta_2}, \pheight{\delta_3}) + 1$۔
\end{enumerate}
که سېکوېنټ~$S$ د $\le n$ لوړوالي !!a{proof} ولري،
$\Proves[n] S$ ليکو۔
\end{defn}

\begin{ex}
په \Log{N1i} کښې هر !!{formula} له خپل ځانه د پرانيستي فرض
په توګه !!a{proof} ګڼل کېږي۔ نو
\[
!A \land !B^x
\]
!!a{proof}~$\delta_1$ دے۔ لوړوالے ئې~$0$ دے۔ له
$\delta_1$ څخه د~\Elim{\land} د دوو ډولونو په کارولو دوه
!!{proof}s~$\delta_2$ او~$\delta_3$ ترلاسه کولے شو:
\[
\AxiomC{$!A \land !B^x$}
\RightLabel{\Elim{\land}[2]}
\UnaryInfC{$!B$}
\DisplayProof
\qquad
\AxiomC{$!A \land !B^x$}
\RightLabel{\Elim{\land}[1]}
\UnaryInfC{$!A$}
\DisplayProof
\]
په~$\delta_2$ او~$\delta_3$ کښې د وروستي استنتاج سمدستي
فرعي !!{proof}s هماغه يو دي: $!A \land !B$۔ دواړه
!!{proof}s د $!A\land !B$ راتګ د پرانيستي فرض په توګه لري۔
$\pheight{\delta_2} = \pheight{\delta_3} = 1$ لرو۔
\textbf{سرچينه‌يي سمون OLCMP-143:} د فرض يوازېني ثبوت
لوړوالے صفر دے؛ د دواړو يو-استنتاجي ثبوتونو لوړوالے يو دے۔
په اصلي برابرۍ کښې د لومړي ثبوت ناسمه شمېره سمه شوه۔

د \Intro{\land} قاعده د~$!B \land !A$ په بڼه پايلې
دپاره د $!B$ او $!A$ په بڼه دوه مقدمې غواړي۔ نو دا
!!a{proof}~$\delta_4$ دے:
\[
\AxiomC{$!A \land !B^x$}
\RightLabel{\Elim{\land}[2]}
\UnaryInfC{$!B$}
\LeftSubproofLabel{$\delta_2$}
\AxiomC{$!A \land !B^x$}
\RightLabel{\Elim{\land}[1]}
\UnaryInfC{$!A$}
\RightSubproofLabel{$\delta_3$}
\RightLabel{\Intro{\land}}
\BinaryInfC{$!B \land !A$}
\DisplayProof
\]
لوړوالے ئې $\pheight{\delta_4} = \max(\pheight{\delta_2},
\pheight{\delta_3})+1 = \max(1,1) + 1 = 2$ دے۔ د فرضونو
په توګه د $!A \land !B$ دوه راتګونه لري او دواړه پرانيستي دي۔

په پاے کښې د \Intro{\lif} په کارولو $\delta_5$ ترلاسه
کولے شو:
\[
\AxiomC{$\Discharge{!A \land !B}{x}$}
\RightLabel{\Elim{\land}[2]}
\UnaryInfC{$!B$}
\LeftSubproofLabel{$\delta_2$}
\AxiomC{$\Discharge{!A \land !B}{x}$}
\RightLabel{\Elim{\land}[1]}
\UnaryInfC{$!A$}
\RightSubproofLabel{$\delta_3$}
\RightLabel{\Intro{\land}}
\BinaryInfC{$!B \land !A$}
\DischargeRule{\Intro{\lif}}{x}
\UnaryInfC{$(!A \land !B) \lif (!B \land !A)$}
\DisplayProof
\]
لوړوالے ئې $\pheight{\delta_4} + 1 = 3$ دے۔ د
\Intro{\lif} استنتاج د $!A \land !B$ په بڼه ټول په~$x$
نښه شوي فرضونه باسي، يعنې د سمدستي فرعي ثبوت دواړه پخوا
پرانيستي فرضونه۔ دا چې د $\delta_5$ يوازې همدغه فرضونه دي،
نو پرانيستے فرض نۀ لري، او د خپل پاييز !!{formula}
$(!A \land !B) \lif (!B \land !A)$ \emph{ثبوت} دے، يعنې
!!a{proof} دے۔
\end{ex}

\end{document}
